\chapter{Mixture of Experts}\label{ch:moe}

\begin{quote}\itshape
One substitution turns Chapter~\ref{ch:em}'s mixture into the thesis
model: let the parameters be functions of the observables. The weights
become a gate evaluated on the encoded market state; the means become
expert networks evaluated on the characteristics. This chapter derives
what that substitution does to training --- above all the gate-logit
gradient $\pi - r$, the two-symbol formula that repairs the prototype's
central defect and that you should be able to reproduce on a napkin for
the rest of your career.
\end{quote}

\section{From mixture to conditional mixture}\label{sec:conditionalmoe}

Chapter~\ref{ch:em} fit one global mixture: every observation faced the
same weights $\pi_k$ and the same component means $\mu_k$. The thesis
model (Definition~\ref{def:nec}) makes both \emph{conditional}:
\begin{equation}\label{eq:moedensity}
  p(y_{i,t} \mid \vx_{i,t}, \vh_t)
  \;=\; \sum_{k=1}^{K} \pi_k(\vh_t)\;
  \gauss\!\bigl(y_{i,t};\, \mu_k(\vx_{i,t}),\, \sigma_k^2\bigr),
\end{equation}
with the three modules of Chapter~\ref{ch:regimes}'s blueprint:
\begin{itemize}[leftmargin=1.6em, itemsep=1pt]
\item the \emph{gate} produces logits $\vz = g(\vh_t) \in \R^K$ from the
  encoded market state and weights $\pi = \operatorname{softmax}(\vz)$
  --- the regime \emph{prior} before seeing the outcome;
\item the \emph{experts} produce means $\mu_k(\vx_{i,t})$ from the
  name-level characteristics --- $K$ regime-conditional prediction
  functions;
\item the \emph{scales} $\sigma_k$ are per-expert free parameters
  (as $\log\sigma_k$), not functions --- each regime carries its own
  noise level, which Chapter~\ref{ch:regimes} argued is most of what a
  regime \emph{is}.
\end{itemize}
Everything in Chapter~\ref{ch:em} survives the substitution wholesale,
because none of its derivations used the fact that $\pi_k$ and $\mu_k$
were free scalars: responsibilities are still \eqref{eq:resp} with
$\pi_k(\vh_t)$ and $\mu_k(\vx_{i,t})$ inserted; the ELBO decomposition,
Fisher's identity, and the GEM interpretation hold verbatim. What changes
is \emph{where gradients go next}: after the chain rule reaches
$\partial/\partial \pi_k$ and $\partial/\partial \mu_k$, it continues into
network weights instead of stopping at scalars. The new mathematics of
this chapter is entirely about that continuation.

\section{The objective, fused}\label{sec:fusedobjective}

Define per-sample joint scores and the loss exactly as the package
computes them:
\begin{equation}\label{eq:fused}
  a_k \;=\; \log \pi_k \;+\; \log \gauss(y; \mu_k, \sigma_k^2),
  \qquad
  \NLL \;=\; -\LSE_k(a_k) \;=\; -\log \sum_k e^{a_k},
\end{equation}
with $\log \pi = \operatorname{log\_softmax}(\vz)$ and the Gaussian
log-density expanded in the log domain,
\[
  \log \gauss(y; \mu_k, \sigma_k^2)
  = -\tfrac12 \log 2\pi - \log \sigma_k
    - \frac{(y - \mu_k)^2}{2\sigma_k^2}.
\]
Numerically, $\LSE$ subtracts the row maximum before exponentiating, so
the pipeline never materializes a density --- only densities' logs ---
and never underflows. The responsibilities emerge from the same tensor:
$\log r_k = a_k - \LSE(a)$, i.e.\ $r = \operatorname{softmax}(a)$. One
fused computation yields objective and posterior; Chapter~\ref{ch:em}'s
E-step is literally a by-product of evaluating the loss.

\begin{whatbreaks}[the prototype's separated pipeline]
The prototype (defect 3) placed \code{nn.Softmax} \emph{inside} the model
and attached a separate cross-entropy loss to the gate. Two distinct
failures follow. Numerical: probabilities near 0 or 1 followed by a
downstream $\log$ is catastrophic cancellation by design ---
$\log(\operatorname{softmax})$ loses precision exactly where confident
routing lives; the fused $\operatorname{log\_softmax}$/$\LSE$ path is
algebraically identical and numerically stable everywhere. Statistical:
a gate trained by its own auxiliary loss against external labels is no
longer maximizing \eqref{eq:moedensity}'s likelihood --- the mixture's
components stop negotiating through a shared objective, and the ``prior''
stops being the quantity the likelihood needs it to be. One model, one
objective, one log domain --- the entire fix.
\end{whatbreaks}

\section{The gate gradient}\label{sec:gategradient}

Here is the book's most important derivation. It is four lines.

\begin{theorem}[The gate learns $\pi - r$]\label{thm:pir}
For the fused objective \eqref{eq:fused} with
$\pi = \operatorname{softmax}(\vz)$,
\begin{equation}\label{eq:pir}
  \frac{\partial\, \NLL}{\partial z_k} \;=\; \pi_k - r_k .
\end{equation}
\end{theorem}

\begin{proof}
Two softmax Jacobian passes, chained. First, $\NLL = -\LSE(a)$ gives
$\partial \NLL / \partial a_j = -\operatorname{softmax}(a)_j = -r_j$.
Second, $a_j = z_j - \LSE(\vz) + \log \gauss_j$ (writing
$\log\pi_j = z_j - \LSE(\vz)$), so
$\partial a_j / \partial z_k = \delta_{jk} - \pi_k$, which does not
depend on the likelihood term. Chain and use $\sum_j r_j = 1$:
\[
  \frac{\partial \NLL}{\partial z_k}
  = \sum_j (-r_j)\,(\delta_{jk} - \pi_k)
  = -r_k + \pi_k \sum_j r_j
  = \pi_k - r_k. \qedhere
\]
\end{proof}

Every property of this formula deserves a sentence.

\emph{It is a Bayes update, imposed as a gradient.} The gate's output
$\pi$ is the prior over regimes; $r$ is the posterior after the outcome
$y$ is seen. Gradient descent moves the logits to shrink $\pi - r$: the
gate is perpetually trained toward \emph{its own posterior}. The
mixture teaches its prior to predict what its likelihood will conclude
--- self-distillation of Bayes' rule, with no labels anywhere.
Decision~B (the gate stays unsupervised) is not an abstinence; it is the
recognition that \eqref{eq:pir} already supplies the supervision, and
that any external label would overwrite the only training signal that
makes ``do the learned regimes mean anything?'' an open, testable
question.

\emph{It is softmax cross-entropy with soft targets.} Classification
training gives logit gradients $p - \bm{e}_{y}$ (predicted probabilities
minus the one-hot label). Formula \eqref{eq:pir} is the same expression
with the hard label replaced by the responsibility vector: the experts
jointly manufacture the label the gate is trained against, each date, each
name. Everything known about cross-entropy training --- calibration at
the optimum, gradient boundedness, saturation behavior --- transfers.

\emph{It vanishes exactly at calibration.} $\partial \NLL / \partial
\vz = 0$ iff $\pi = r$: the stationary gate is the one whose prior
already matches the posterior on average
(Exercise~\ref{ex:5-calibration} makes the ``on average'' precise). A
well-trained gate is thus a \emph{calibrated} regime predictor by
construction --- the property Chapter~18's reliability diagnostics test
empirically.

\emph{It is what the prototype destroyed.} Route through
$\operatorname{argmax}$ instead of the mixture and the entire derivation
collapses at line one: the loss ceases to depend on $\vz$ except through
a piecewise-constant selection, and the gradient is zero almost
everywhere (Chapter~\ref{ch:autopsy}, defect 1;
Chapter~\ref{ch:routing} proves the subtler renormalized variant). The
package's regression test \code{test\_gate\_receives\_gradient} asserts
\eqref{eq:pir} is alive after one backward pass --- and its companion
asserts the argmax double is exactly dead. The two-line formula is
guarded by a two-assertion test.

\section{The expert and scale gradients}\label{sec:expertgradients}

Differentiating \eqref{eq:fused} with respect to the remaining outputs
(Exercise~\ref{ex:5-verify} checks all of these against autograd):
\begin{align}
  \frac{\partial\, \NLL}{\partial \mu_k}
  &= r_k\, \frac{\mu_k - y}{\sigma_k^2},
  \label{eq:mugrad}\\[2pt]
  \frac{\partial\, \NLL}{\partial \log \sigma_k}
  &= r_k \left(1 - \frac{(y - \mu_k)^2}{\sigma_k^2}\right).
  \label{eq:siggrad}
\end{align}
Equation \eqref{eq:mugrad} is \emph{responsibility-weighted regression}:
expert $k$ receives an ordinary squared-error pull toward $y$, scaled by
how responsible the posterior holds it. Experts with $r_k \approx 0$ on a
sample are untouched by it --- specialization is not just permitted but
enforced by the gradient's structure. Equation \eqref{eq:siggrad} is
stationary when $\sigma_k^2$ equals the responsibility-weighted mean
squared residual \emph{of expert $k$'s own samples}: each regime's noise
level is estimated from the data that regime claims. Summing over a batch
and setting to zero recovers exactly Chapter~\ref{ch:em}'s M-step
\eqref{eq:mstep} --- as it must, by Fisher's identity.

The full division of labor, in one table:

\begin{center}
\small
\begin{tabular}{@{}llll@{}}
\toprule
quantity & gradient & EM reading & failure if broken \\
\midrule
gate logits $\vz$ & $\pi - r$ & prior $\to$ posterior
  & defect 1 (dead gate) \\
means $\mu_k$ & $r_k(\mu_k - y)/\sigma_k^2$ & weighted M-step
  & pooled experts (Ch.~\ref{ch:symmetry}) \\
scales $\log\sigma_k$ & $r_k\bigl(1 - (y-\mu_k)^2/\sigma_k^2\bigr)$
  & weighted variance & spikes / flat $r$ (Ch.~\ref{ch:em}) \\
responsibilities $r$ & (not a parameter) & E-step, in the forward pass
  & --- \\
\bottomrule
\end{tabular}
\end{center}

\begin{origins}[Adaptive mixtures of local experts]
The architecture and the competitive-training insight are from Jacobs,
Jordan, Nowlan and Hinton, ``Adaptive Mixtures of Local Experts''
(\emph{Neural Computation}, 1991). Their central observation is worth
retelling because it is a loss-design lesson. An earlier ``committee''
objective, $\bigl(y - \sum_k \pi_k \mu_k\bigr)^2$, trains the
\emph{blended} prediction: each expert's target is the residual left by
the blend, every expert ends up correcting every other, and the ensemble
\emph{cooperates} into homogeneity. Switching to the mixture
(log-)likelihood --- $\sum_k \pi_k \exp(-\tfrac12(y-\mu_k)^2)$ in their
Gaussian case --- makes experts \emph{compete} for responsibility on each
case, and specialization follows. The moral, which recurs in
Chapter~\ref{ch:routing}: with mixtures, the difference between an
architecture that specializes and one that homogenizes is not the wiring
but the objective. Jordan and Jacobs (1994) supplied the EM view and
hierarchical gates; Shazeer et al.\ (2017) scaled sparse variants to
language models, whose load-balancing pathologies reappear in our
Section~\ref{sec:routing-family} context. The thesis model is the 1991
design, taken seriously with 2020s hygiene.
\end{origins}

\section{Prediction and uncertainty}\label{sec:prediction}

Training maximizes a density; the desk needs numbers. The model's point
prediction is the mixture mean,
\begin{equation}\label{eq:yhat}
  \hat y \;=\; \E[y \mid \vx, \vh]
  \;=\; \sum_k \pi_k\, \mu_k(\vx),
\end{equation}
--- the posterior-weighted blend, \emph{not} the best expert's output.
Selecting $\mu_{\argmax_k \pi_k}$ discards the model's own uncertainty
about the regime and reintroduces a discontinuity (as the gate wavers
across $0.5$, predictions jump), for zero statistical benefit: under
squared error the conditional mean is optimal, and under rank-IC grading
any monotone transform of it scores identically.

The predictive \emph{variance} decomposes by the law of total variance
(Exercise~\ref{ex:5-variance}):
\begin{equation}\label{eq:totalvar}
  \Var(y \mid \vx, \vh)
  = \underbrace{\sum_k \pi_k \sigma_k^2}_{\text{within-regime noise}}
  \;+\;
  \underbrace{\sum_k \pi_k\, \mu_k^2 - \hat y^{\,2}}_{
    \text{between-regime disagreement}} .
\end{equation}
The second term is zero exactly when the experts agree; it is the model's
self-reported ``the regimes disagree about this name'' signal ---
information a single-Gaussian model structurally cannot emit, and one of
the concrete payoffs of carrying a full conditional density rather than a
point estimate (Section~\ref{sec:whatisresult}'s argument, now with a
formula).

\section{The interface: log-prior versus training weights}
\label{sec:interface}

One design decision in the package's loss layer belongs to the theory and
must be stated now. The fused objective \eqref{eq:fused} consumes
\emph{some} log-weight vector alongside the expert log-likelihoods. For
the soft gate the choice is unambiguous: the log-prior
$\operatorname{log\_softmax}(\vz)$, both for training and for held-out
evaluation. Chapter~\ref{ch:routing}'s hard-routing variant will need the
two roles to \emph{differ}: the quantity that makes its objective
trainable (a hard-EM surrogate) is not the quantity that honestly states
its predictive distribution (the selected expert's density). The
package's \code{PriorOutput} therefore carries \code{log\_prior} always,
and optionally \code{log\_train\_weights}; the trainer uses the training
weights when present \emph{for the loss only}, and every held-out NLL in
every table of Part~V is computed under \code{log\_prior} --- the model's
honest predictive stance. This is a small interface with a large moral:
\emph{when a training surrogate and a predictive distribution diverge,
name both, and never grade on the surrogate.}

\begin{codewalk}[nec\_moe/model.py + likelihood.py --- assembly]
The forward pass assembles \eqref{eq:moedensity} in five steps, each one
module, each module owning exactly one formula of this chapter:
\begin{lstlisting}
h_t = self.encoder(x_seq)                  # (B, n)   market state
z   = self.gate(h_t)                       # (B, K)   logits -- no softmax!
out = self.prior(z, ctx)                   # log pi   (pluggable, Ch. 6-7)
mu, log_sigma = self.experts(x_exp)        # (B, K), (K,)
log_lik = expert_log_likelihood(mu, log_sigma, y)     # (B, K)
nll_out = mixture_nll(out.log_prior, log_lik)         # fused (Ch. 4)
\end{lstlisting}
The gate module ends at logits --- the softmax lives inside the fused
loss (defect-3 fix, enforced by the architecture rather than by
discipline). \code{expert\_log\_likelihood} is a pure, parameter-free
function computing the Gaussian log-density directly in log domain; its
purity is what makes ``same likelihood, different prior'' literally true
when Chapter~\ref{ch:hmm} swaps the prior box. And the un-fused
alternative --- \code{prob.log()} after a softmax, densities
exponentiated then logged --- is not merely discouraged: no code path in
the package can express it.
\end{codewalk}

\exercisehead

\begin{exercise}\label{ex:5-verify}
(Hands-on.) Build a tiny mixture forward pass in raw torch (no package):
random $\vz \in \R^{1\times 3}$ with \code{requires\_grad}, fixed
$\mu \in \R^{1 \times 3}$, $\sigma = 1$, scalar $y$. Compute the fused
NLL and \code{backward()}. Verify numerically that \code{z.grad} equals
$\pi - r$, that \code{mu.grad} equals \eqref{eq:mugrad}, and (promoting
$\log\sigma$ to a parameter) that its gradient matches
\eqref{eq:siggrad}. Then break it: recompute with
\code{torch.where(argmax(z)...)} routing and confirm \code{z.grad} is
zero --- you have reproduced both halves of the package's defect-1
regression test.
\end{exercise}

\begin{exercise}\label{ex:5-sumzero}
Show from \eqref{eq:pir} that $\sum_k \partial \NLL/\partial z_k = 0$
for every sample. Connect this to the softmax's invariance under adding a
constant to all logits, and explain why the gate's logits therefore never
receive a signal about their overall level --- only about their
differences. What (harmless) consequence does this have for the learned
logits over a long training run, and how does the gate head's
\code{LayerNorm} + weight decay (Chapter~9) keep it tidy?
\end{exercise}

\begin{exercise}\label{ex:5-variance}
Derive \eqref{eq:totalvar} from the law of total variance conditioning
on the component index. Then compute both terms for a two-expert model
with $\pi = (0.5, 0.5)$, $\mu = (+1, -1)$, $\sigma = (1, 1)$ and for
$\pi = (0.95, 0.05)$, same $\mu, \sigma$. Which situation would a risk
manager rather know about, and which summary --- $\hat y$ alone or
\eqref{eq:totalvar}'s split --- reveals it?
\end{exercise}

\begin{exercise}\label{ex:5-calibration}
Define the gate's \emph{calibration residual} as
$\E\bigl[r \mid \vh\bigr] - \pi(\vh)$. Show that the expected gate-logit
gradient at fixed $\vh$ is exactly its negative, so a stationary gate
satisfies $\pi(\vh) = \E[r \mid \vh]$: the prior equals the average
posterior, conditionally. Explain (two sentences) why this is precisely
the property that Chapter~18's reliability diagram tests on held-out
data, and why finding it \emph{violated} out-of-sample is evidence of
regime drift rather than a training bug.
\end{exercise}

\begin{exercise}\label{ex:5-decisionA}
(Decision A's gradient anatomy.) Suppose the experts read the
concatenation $[\vx; \vh_t]$ (\code{input\_mode="snapshot\_plus\_hidden"})
instead of $\vx$ alone. Write the gradient of the per-sample NLL with
respect to the encoder state $\vh_t$ as a sum of a gate path and an
expert path, and show the expert path is
$\sum_k r_k \frac{\mu_k - y}{\sigma_k^2}\,
\partial \mu_k / \partial \vh_t$. In the prototype, defect 1 killed the
gate path and defect 6 severed the expert path --- conclude what total
gradient its four-million-parameter encoder received from the prediction
task, and why the corrected default (snapshot-only experts) makes the
gate path the encoder's \emph{sole}, and therefore interpretable,
training signal.
\end{exercise}

\begin{exercise}\label{ex:5-jjnh}
Reproduce the 1991 cooperation/competition contrast in numpy: two linear
experts, a fixed uniform gate, data from the sign-flip world
(Proposition~\ref{prop:signflip}). Train once on the committee loss
$\bigl(y - \tfrac12(\mu_1 + \mu_2)\bigr)^2$ and once on the mixture NLL,
from the same asymmetric initialization $(b_1, b_2) = (0.8, -0.6)$.
Report the terminal slopes. One objective drives both experts to the
pooled solution; the other preserves specialization. Explain each
outcome from the two losses' gradients.
\end{exercise}
